% =====================================================================
%  FULL PAPER -- EMPIRICAL / SYSTEMS VERSION (Option B)
%  Elsevier elsarticle / JPDC
%
%  "Online Capacitated Routing for Mixture-of-Experts:
%   A Single-Pass Water-Filling Router for Load Balancing
%   without Auxiliary Losses"
%
%  WHAT CHANGED vs the earlier draft (and WHY):
%   - REMOVED the competitive-ratio theorem and the max-load theorem.
%     A verified counterexample shows the (1-1/e) claim is FALSE for
%     count-capacity, continuous gate scores under adversarial order;
%     and WFPD's monotone price does NOT reach near-optimality under
%     random order either (verified numerically: ALG/OPT plateaus ~0.85,
%     not -> 1). A provably (1-eps) variant needs a different price
%     update (dual mirror descent) -- that is a separate paper.
%   - This version is an honest empirical/systems contribution: online
%     packing FORMULATION + LP optimum as an EVALUATION BENCHMARK +
%     a lightweight water-filling heuristic + constant-time complexity
%     (the one theorem that is true) + strong trace-driven experiments.
%
%  Compile: pdflatex (bibliography is inline). Needs: elsarticle,
%  amsthm, algorithm, algpseudocode, booktabs. Times via newtx
%  (delete the two newtx lines if your install clashes).
%
%  >>> VERIFY all \bibitem details before submission.
% =====================================================================
\documentclass[preprint,12pt]{elsarticle}

\usepackage{newtxtext,newtxmath}
\usepackage{amsmath,amssymb,amsthm}
\usepackage{algorithm}
\usepackage{algpseudocode}
\usepackage{booktabs}
\usepackage{url}

\theoremstyle{plain}
\newtheorem{theorem}{Theorem}
\theoremstyle{remark}
\newtheorem{remark}{Remark}

\newcommand{\opt}{\mathrm{OPT}}
\newcommand{\alg}{\mathrm{ALG}}
\newcommand{\Lmax}{L_{\max}}

\journal{Journal of Parallel and Distributed Computing}

\begin{document}

\begin{frontmatter}

\title{Online Capacitated Routing for Mixture-of-Experts:
A Single-Pass Water-Filling Router for Load Balancing without Auxiliary
Losses}

\author{Kefan Liu}
\ead{kefan001@e.ntu.edu.sg}
\affiliation{organization={Centre in Computational Technologies for Finance (CCTF), Nanyang Technological University},
             city={Singapore},
             country={Singapore}}

\begin{abstract}
Mixture-of-Experts (MoE) models route each token to a few of many experts,
and balanced expert load is essential for parallel utilisation: the
throughput of an MoE layer is set by its busiest expert. Yet the standard
top-$k$ gate balances poorly and is propped up by auxiliary load-balancing
losses and capacity-dropping heuristics that are tied to training and
require tuning. We formulate inference-time expert routing as an
\emph{online capacitated packing} problem---tokens stream in and each is
routed to at most $k$ of $m$ capacity-limited experts to maximise total
gate score---and use its offline linear-programming optimum as a strong
benchmark for measuring routing quality. We propose WFPD, a lightweight
single-pass \emph{water-filling} router: each expert carries a price that
rises as it fills, and each token is sent to its top-$k$ positive-surplus
experts, so load balancing emerges from the routing rule itself, with no
auxiliary loss, no token-dropping heuristic and no tuning. WFPD updates in
$O(m+k\log C)$ time and $O(m)$ space per token, independent of sequence
length. On real routing traces from \textsc{OLMoE-1B-7B}, WFPD consistently
improves over the standard token-choice gate, retaining $95\%$ of the
offline optimum versus $88\%$ while halving the load imbalance, and matches
an iterative optimal-transport solver while balancing better---in a single
online pass at stable throughput.
\end{abstract}

\begin{keyword}
mixture-of-experts \sep online routing \sep load balancing \sep
resource allocation \sep inference systems
\end{keyword}

\end{frontmatter}

% =====================================================================
\section{Introduction}
\label{sec:intro}

Mixture-of-Experts (MoE) layers scale neural networks without a
proportional increase in compute: each token is processed by only a few of
many expert sub-networks, so the parameter count can grow while the
per-token cost stays fixed~\cite{shazeer2017,switch}. A lightweight
\emph{router} decides, for every token, which $k$ of the $m$ experts handle
it. Because the experts run in parallel, often sharded across devices, the
throughput of an MoE layer is set by its \emph{busiest} expert: if the
router concentrates tokens on a few popular experts the rest sit idle,
capacity is wasted, and overloaded experts drop tokens, hurting
quality~\cite{gshard,switch}. Keeping the per-expert load balanced is thus a
parallel resource-allocation problem at the heart of MoE efficiency.

The prevailing remedies are heuristic and training-coupled. The dominant
\emph{token-choice} gate---each token greedily takes its $k$ highest-scoring
experts---is balanced only by adding an \emph{auxiliary load-balancing
loss} to the training objective and by a \emph{capacity factor} that drops
tokens once an expert is full~\cite{switch,gshard}. These devices tangle
balancing into training, need tuning, and still leave the realised load
uneven. \emph{Expert-choice} routing~\cite{expertchoice} inverts the
assignment so each expert selects its favourite tokens, which balances load
but assigns a variable number of experts per token and needs the whole
batch at once. \textsc{BASE} layers~\cite{baselayers} solve a global
linear-assignment problem per layer, again requiring all tokens
simultaneously.

We take an online view. At inference time, routing is fundamentally an
online problem: tokens arrive one at a time, the router must commit each
before seeing the next, and every expert has a finite capacity. We
formulate inference-time MoE routing as \emph{online capacitated packing}.
This formulation does two things: it gives a clean, assumption-light
statement of the routing task, and it yields the offline LP optimum as a
strong yardstick against which any router can be measured. We then design a
simple router for this setting and evaluate it on real MoE routing traces.

\paragraph{Contributions.}
\begin{enumerate}
\item We formulate inference-time MoE routing as an online capacitated
packing problem (Section~\ref{sec:problem}) and use its offline linear-%
programming optimum as a benchmark for routing quality.
\item We propose \emph{WFPD}, a lightweight single-pass water-filling router
(Section~\ref{sec:algo}). A per-expert price rises as the expert fills and
each token is routed to its top-$k$ positive-surplus experts, so load
balancing emerges from the routing rule---no auxiliary loss, no token
dropping, no tuning---with $O(m+k\log C)$ time and $O(m)$ space per token,
independent of sequence length (Theorem~\ref{thm:complexity}).
\item We evaluate WFPD on real routing traces from
\textsc{OLMoE-1B-7B} (Section~\ref{sec:exp}). It consistently improves on
the standard token-choice gate in both retained gate score and load
balance, and matches an iterative optimal-transport solver while balancing
better, in a single online pass at stable throughput.
\end{enumerate}

We keep the contribution empirical and systems-oriented. As we discuss in
Section~\ref{sec:algo}, no fixed online rule can guarantee a constant
fraction of the offline optimum for continuous gate scores under
adversarial arrival order; designing a router with provable guarantees in a
stochastic or random-order model is a separate question that we leave to
future work (Section~\ref{sec:conclusion}).

% =====================================================================
\section{Related Work}
\label{sec:related}

\paragraph{Routing in Mixture-of-Experts.}
Sparsely-gated MoE layers were introduced by Shazeer et
al.~\cite{shazeer2017} and scaled by GShard~\cite{gshard} and Switch
Transformer~\cite{switch}, which established the now-standard recipe of a
token-choice gate regularised by an auxiliary load-balancing loss and a
capacity factor with token dropping. Recent open models such as
Mixtral~\cite{mixtral} and \textsc{OLMoE}~\cite{olmoe} follow the same
template. Two lines change the assignment to improve balance:
Expert-Choice routing~\cite{expertchoice} lets each expert pick its top
tokens, and \textsc{BASE} layers~\cite{baselayers} solve a global linear
assignment per layer; both are full-batch and trade away the streaming,
per-token structure of autoregressive inference. WFPD keeps the
token-choice interface---one token at a time, at most $k$ experts each---but
replaces the greedy gate with a priced one that balances load
intrinsically, without an auxiliary loss.

\paragraph{Online resource allocation.}
Our formulation places MoE routing in the setting of online capacitated
assignment, which has a rich theory: online bipartite
matching~\cite{kvv}, budgeted allocation and the AdWords
problem~\cite{adwords}, and the online primal--dual
framework~\cite{buchbinder}. In the random-order and stochastic models,
online packing LPs admit near-optimal guarantees once capacities are
large~\cite{devanurhayes,awy,krtv,blm}. We use this lens for problem
framing and for benchmarking against the offline optimum; designing a
router with provable competitive guarantees in this model---which, as we
note in Section~\ref{sec:algo}, requires a price update different from
ours---is a direction we leave to future work.

\paragraph{Load balancing.}
The balls-into-bins and power-of-two-choices
literature~\cite{twochoices,mitzenmacher} studies how to keep the maximum
load close to the mean; WFPD's price mechanism makes a token less likely to
land on a heavily-loaded expert, in the same spirit, but as a deterministic
routing rule rather than a randomised allocation.

\paragraph{Optimal transport for balanced assignment.}
Balanced routing can be cast as an optimal-transport problem and solved with
Sinkhorn iterations~\cite{sinkhorn}; we use such a solver (\textsc{LPR}) as
a strong quality baseline. It is iterative and full-batch, with an iteration
count and a convergence tolerance to tune, whereas WFPD operates in a single
online pass.

% =====================================================================
\section{Problem Formulation}
\label{sec:problem}

\paragraph{Setup.}
A stream of $n$ tokens arrives one at a time and must be routed among $m$
experts. Token $t$ is revealed with its \emph{gate scores}
$g_{t,e}\in[0,1]$ for every expert $e\in[m]$; in practice $g_{t,e}$ is the
softmax of the router logit. Each token may be sent to at most $k$ experts
(the top-$k$ budget), and each expert can process at most $C$ tokens (its
\emph{capacity}). The objective is to maximise the total served gate score
$\alg=\sum_{t}\sum_{e\in S_t} g_{t,e}$, where $S_t$ is the set of experts
chosen for token $t$ and $|S_t|\le k$. In the online model, $S_t$ is
committed before token $t+1$ is revealed and is irrevocable, and an expert
that has accepted $C$ tokens is unavailable.

\paragraph{Offline optimum as a benchmark.}
To measure how much routing quality a router achieves, we compare its served
gate score against the clairvoyant offline optimum of the capacitated
assignment LP, with $x_{t,e}\in[0,1]$:
\begin{equation}
\label{eq:primal}
\begin{aligned}
  \opt=\max_{x\ge 0}\ & \sum_{t,e} g_{t,e}\,x_{t,e}\\
  \text{s.t.}\quad
  & \textstyle\sum_{e} x_{t,e}\le k && \forall t\quad(\text{budget})\\
  & \textstyle\sum_{t} x_{t,e}\le C && \forall e\quad(\text{capacity})\\
  & x_{t,e}\le 1 && \forall t,e.
\end{aligned}
\end{equation}
$\opt$ is the largest total gate score attainable by any assignment that
sends each token to at most $k$ experts and never exceeds capacity $C$. It
is not attainable online---it requires seeing all tokens at once---but it is
a tight upper bound and serves as our evaluation yardstick: we report each
router's served score as a fraction of $\opt$ (the \emph{retained} metric of
Section~\ref{sec:exp}).

% =====================================================================
\section{The WFPD Router}
\label{sec:algo}

\paragraph{Design rationale.}
WFPD attaches to each expert a \emph{price} $\beta_e$---a congestion signal,
or water level, that rises as the expert fills. A token sees the
\emph{surplus} $\sigma_{t,e}=g_{t,e}-\beta_e$ of each expert and is routed to
the (at most $k$) experts of largest positive surplus. Whenever an expert
serves a token its price is raised multiplicatively, so popular experts
become progressively more expensive and later tokens are deflected toward
cheaper, less-loaded alternatives. Load balancing is thus emergent from the
routing rule: there is no auxiliary load-balancing loss, no
capacity-dropping heuristic and no global view of the layer. The growth rate
is calibrated so that the price approaches saturation as the expert
approaches capacity. WFPD is summarised in Algorithm~\ref{alg:wfpd}.

\begin{algorithm}[t]
\caption{Water-Filling Router (WFPD)}
\label{alg:wfpd}
\begin{algorithmic}[1]
\Require experts $[m]$, capacity $C$, budget $k$; stream of gate scores
\State $\beta_e\gets 0,\ \ n_e\gets 0$ \textbf{ for all } $e\in[m]$
   \Comment{prices and per-expert load counters}
\State $M\gets (1+1/C)^{C}-1$
\For{each arriving token $t$ with scores $(g_{t,e})_{e\in[m]}$}
  \State $\sigma_{t,e}\gets g_{t,e}-\beta_e$ \textbf{ for all } $e$ with $n_e<C$
         \Comment{surplus of admissible experts}
  \State $S_t\gets$ the (at most $k$) experts with largest $\sigma_{t,e}>0$
  \For{$e\in S_t$}
     \State route token $t$ to expert $e$;\ \ $n_e\gets n_e+1$
     \State $\beta_e\gets \beta_e\,(1+1/C)+g_{t,e}/(C\,M)$
            \Comment{raise the water level}
  \EndFor
\EndFor
\end{algorithmic}
\end{algorithm}

\paragraph{Scope.}
WFPD is a heuristic, and we are deliberate about what we claim for it. The
price update is an online analogue of water-filling, but it does \emph{not}
come with a worst-case competitive guarantee. Under adversarially ordered
inputs with continuous gate scores, no fixed online routing rule can
guarantee a constant fraction of the offline optimum: a low-score token can
occupy a slot that a later high-score token would have used, and an
adversary can force this on any algorithm that must commit irrevocably. Our
design instead targets realistic token streams, on which---as
Section~\ref{sec:exp} shows---the rule performs strongly. A variant with
provable guarantees under random-order or stochastic arrivals would require
a price update that can both rise and fall toward a target utilisation
(rather than the monotone update above); we regard this as a separate
contribution and leave it to future work.

\begin{theorem}[Per-token complexity]
\label{thm:complexity}
WFPD processes each token in $O(m+k\log C)$ time and uses $O(m)$ space.
\end{theorem}

\begin{proof}
The state consists of the prices $\beta_e$ and counters $n_e$, i.e.\ $O(m)$
space. Per token, evaluating the $m$ surpluses costs $O(m)$. Selecting the
$k$ largest positive surpluses and updating the prices and counters of those
$k$ experts costs $O(k\log C)$ using a priority queue keyed by surplus, the
$\log C$ factor accounting for the reorganisation of an expert as its load
counter advances through its $C$ capacity levels. The total $O(m+k\log C)$
is independent of the sequence length $n$.
\end{proof}

\begin{remark}
Theorem~\ref{thm:complexity} is the formal counterpart of the empirical
observation (Section~\ref{sec:exp-throughput}) that WFPD's per-token cost is
flat in $n$: the price bookkeeping adds a constant overhead per token, never
a growing one, which makes the router suitable for streaming, single-pass
inference.
\end{remark}

% =====================================================================
\section{Experimental Evaluation}
\label{sec:exp}

We evaluate WFPD against three baselines on the routing traces of a real
MoE language model, asking how much of the offline routing quality it
retains, how well it balances expert load, and what its per-token cost is.

\subsection{Setup}
\label{sec:exp-setup}

\paragraph{Model and corpus.}
We use \textsc{OLMoE-1B-7B}~\cite{olmoe} ($6.9$B parameters; $1$B active),
which routes each token to $k=8$ of $m=64$ experts in every one of its $16$
MoE layers. We run the model in $4$-bit precision on a single commodity GPU
and feed it $128$ documents from \textsc{WikiText-2}~\cite{wikitext}. At
each token position we record the router logits of every MoE layer; after a
softmax these give the per-token gate scores $g_{t,e}\in[0,1]$. This
produces roughly $8.8\times10^{3}$ token-to-expert routing events per layer,
aggregated across all layers for the statistics below.

\paragraph{Capacity and the offline optimum.}
For $n$ tokens we set the per-expert capacity to $C=\lceil\gamma\,nk/m\rceil$
with slack $\gamma=1.25$. For each layer we compute the offline optimum
$\opt$ of \eqref{eq:primal} exactly with the HiGHS LP solver and report
quality relative to it.

\paragraph{Metrics.}
\textbf{Retained} $=\alg/\opt$, the fraction of the offline-optimal total
gate score captured online; \textbf{Gini}, the Gini coefficient of the
per-expert load ($0=$ balanced); $\boldsymbol{\Lmax/C}$, the busiest
expert's load relative to capacity; \textbf{Min/Max}, the ratio of the
least- to the most-loaded expert ($0=$ a starved expert); \textbf{Exp./tok}, the average number of experts per token (must be $\le k$
for a budget-respecting router).

\paragraph{Baselines.}
\textsc{Token-choice} is the standard top-$k$ gate: each token greedily
takes its $k$ best experts and is dropped from any whose capacity is full.
\textsc{Expert-choice}~\cite{expertchoice} inverts the assignment---each
expert selects its top-$C$ tokens---balancing load perfectly but not
respecting a per-token budget. \textsc{LPR (OT)} solves an entropic
optimal-transport relaxation with Sinkhorn iterations~\cite{sinkhorn}; it is
iterative and full-batch and serves as a strong offline-style quality
reference.

\subsection{Routing Quality}
\label{sec:exp-quality}

Table~\ref{tab:quality} reports routing quality for the four routers across
all metrics; we discuss each comparison in turn.

\begin{table}[t]
\centering
\caption{Routing quality on \textsc{OLMoE-1B-7B} ($m=64$, $k=8$),
aggregated over all MoE layers ($\approx 8.8$k routings/layer). Higher
\emph{Retained}, \emph{Min/Max} and lower \emph{Gini} are
better. Bold marks the best value among the three \emph{budget-respecting}
routers.}
\label{tab:quality}
\begin{tabular}{lrrrrr}
\toprule
Router & Retained & Gini & $\Lmax/C$ & Min/Max & Exp./tok \\
\midrule
WFPD (ours)               & 0.953          & \textbf{0.106} & 1.00 & \textbf{0.380} & 7.78 \\
Token-choice              & 0.877          & 0.245          & 1.00 & 0.069          & 6.75 \\
Expert-choice$^{\dagger}$ & 1.102          & 0.000          & 1.00 & 1.000          & 10.00 \\
LPR (OT)                  & \textbf{0.969} & 0.141          & 1.00 & 0.267          & 8.00 \\
\bottomrule
\end{tabular}
\\[3pt]
{\footnotesize $\dagger$~\textsc{Expert-choice} does not respect the
per-token budget (Exp./tok $=10>k=8$); reported for reference only.}
\end{table}

\paragraph{WFPD improves on the standard gate on both axes.}
Against \textsc{Token-choice}---the rule deployed in real MoE models---WFPD
is better on \emph{both} axes at once. It retains $95.3\%$ of the offline
optimum versus $87.7\%$, a $7.6$-point ($8.7\%$ relative) gain in served
gate score, while simultaneously more than halving the load imbalance (Gini
$0.106$ vs.\ $0.245$). The gap is sharpest in the tail:
\textsc{Token-choice} drives its least-used expert down to $6.9\%$ of the
busiest expert's load (Min/Max $=0.069$), whereas WFPD holds it at $38.0\%$.
The greedy gate also \emph{wastes} budget---its $6.75$ experts/token, well
below $k=8$, are tokens whose preferred experts were saturated and were
discarded---while WFPD's priced routing redirects those tokens to
near-optimal alternatives, spending $7.78$ experts/token.

\paragraph{Competitive with the OT solver, better balanced.}
\textsc{LPR} retains marginally more gate score than WFPD ($0.969$ vs.\
$0.953$), as expected of an iterative full-batch optimiser with a global
view. Yet WFPD is the better balanced of the two (Gini $0.106$ vs.\ $0.141$;
Min/Max $0.380$ vs.\ $0.267$) and, as Section~\ref{sec:exp-throughput}
shows, achieves this in a single online pass with constant per-token work
and no iteration count or convergence tolerance to tune.

\paragraph{On Expert-choice.}
\textsc{Expert-choice} appears flawless---zero Gini, Min/Max $=1$, Retained
$>1$---but only because it solves a \emph{less constrained} problem: it
assigns a variable number of experts per token ($10$ on average, exceeding
$k=8$), letting it both fill every expert exactly and surpass the top-$k$
optimum that defines $\opt$. It is not a valid competitor in the online
top-$k$ setting and is included only to indicate what relaxing the per-token
budget would buy.

\paragraph{Reading the balance metrics.}
Every router reports $\Lmax/C\approx 1$: at $\gamma=1.25$ slack all methods
saturate the busiest expert, so Gini and Min/Max---not $\Lmax/C$---are the
discriminating balance statistics here. WFPD's $7.78\le k$ experts/token
confirms it respects the per-token budget (unlike Expert-choice), and its
retained value of $0.953$ shows it captures almost all of the offline
optimum on this real workload while keeping load far more even than the
deployed gate.

\subsection{Throughput}
\label{sec:exp-throughput}

We measure per-token routing latency as the number of tokens $n$ grows,
averaged over $15$ trials at $m=64$ (routing only; expert compute
excluded). WFPD costs $\approx 7.3\,\mu$s per token and, crucially, this
cost is flat in $n$ (from $n=256$ to $n=4096$), matching the constant
per-token complexity of Theorem~\ref{thm:complexity}. \textsc{Token-choice}
is slightly cheaper ($\approx 5.3\,\mu$s), \textsc{Expert-choice} is the
fastest ($\approx 0.7\,\mu$s) as it reduces to one per-expert top-$C$
selection, and \textsc{LPR} ranges over $\approx 5$--$7\,\mu$s with variance
that grows in $n$ because of its Sinkhorn iterations. WFPD is not the
wall-clock-fastest router, and speed is not its contribution; its advantage
is structural---online and single-pass, iteration-free, $O(m)$ state, and
free of any auxiliary loss or tuned capacity-dropping heuristic---at a
small, constant per-token overhead.

\subsection{Discussion}
\label{sec:exp-discussion}

On real MoE routing traces, WFPD consistently improves on the standard
\textsc{Token-choice} gate in both retained gate score and load balance,
matches an iterative optimal-transport solver to within $1.6$ points of gate
score while balancing load better than it, and respects the per-token expert
budget---all in a single online pass at constant per-token cost. The
formulation of routing as online capacitated packing, with the offline LP as
a yardstick, makes these comparisons precise; the water-filling rule turns
that yardstick into a practical, parameter-free router.

% =====================================================================
\section{Conclusion}
\label{sec:conclusion}

We formulated inference-time MoE routing as an online capacitated packing
problem and introduced WFPD, a lightweight single-pass water-filling router
that balances expert load as an intrinsic property of its routing rule, with
no auxiliary loss, no token-dropping heuristic and no tuning, at constant
$O(m+k\log C)$ cost per token. On real \textsc{OLMoE-1B-7B} traces it
consistently improves on the standard token-choice gate on both quality and
balance and matches an iterative optimal-transport solver while balancing
better, in a single pass.

Several directions remain. First, the online-packing formulation invites a
router with \emph{provable} guarantees in the random-order or stochastic
arrival model; as we noted, this requires a price update that can both rise
and fall toward a target utilisation---a dual mirror-descent
scheme~\cite{blm}---rather than the monotone update used here, and is a
natural follow-up. Second, our study uses one MoE family; a sweep across
architectures and scales would further test the router. Third, integrating
WFPD as a drop-in inference router and measuring end-to-end throughput on
expert-parallel hardware---where a lower maximum load should translate into
wall-clock speedups---is a systems follow-up well suited to a
parallel-computing setting.

% =====================================================================
%  ELSEVIER-REQUIRED STATEMENTS (after Conclusion, before References)
% =====================================================================
\section*{CRediT authorship contribution statement}
\textbf{Kefan Liu:} Conceptualization, Methodology, Software, Formal
analysis, Investigation, Data curation, Writing -- original draft, Writing
-- review \& editing.

\section*{Declaration of competing interest}
The author declares that she has no known competing financial interests or
personal relationships that could have appeared to influence the work
reported in this paper.

\section*{Funding}
This research did not receive any specific grant from funding agencies in
the public, commercial, or not-for-profit sectors.

\section*{Data availability}
The source code, the routing-trace extraction pipeline and the extracted
\textsc{OLMoE-1B-7B} traces used in this study are publicly available at
\url{https://github.com/liukefan821/moe-online-routing}.

\section*{Declaration of generative AI and AI-assisted technologies in the
manuscript preparation process}
During the preparation of this work the author used Claude (Anthropic) in
order to assist with drafting and editing the manuscript text. After using
this tool the author reviewed and edited the content as needed and takes
full responsibility for the content of the published article.

% =====================================================================
%  BIBLIOGRAPHY -- best-effort; VERIFY every entry before submission.
% =====================================================================
\begin{thebibliography}{99}

\bibitem{shazeer2017} N.~Shazeer, A.~Mirhoseini, K.~Maziarz, A.~Davis,
Q.~Le, G.~Hinton, J.~Dean, Outrageously large neural networks: The
sparsely-gated mixture-of-experts layer, in: ICLR, 2017.

\bibitem{gshard} D.~Lepikhin, H.~Lee, Y.~Xu, D.~Chen, O.~Firat, Y.~Huang,
M.~Krikun, N.~Shazeer, Z.~Chen, {GShard}: Scaling giant models with
conditional computation and automatic sharding, in: ICLR, 2021.

\bibitem{switch} W.~Fedus, B.~Zoph, N.~Shazeer, Switch transformers:
Scaling to trillion parameter models with simple and efficient sparsity,
Journal of Machine Learning Research 23 (120) (2022) 1--39.

\bibitem{expertchoice} Y.~Zhou, T.~Lei, H.~Liu, N.~Du, Y.~Huang, V.~Zhao,
A.~Dai, Z.~Chen, Q.~Le, J.~Laudon, Mixture-of-experts with expert choice
routing, in: NeurIPS, 2022.

\bibitem{baselayers} M.~Lewis, S.~Bhosale, T.~Dettmers, N.~Goyal,
L.~Zettlemoyer, {BASE} layers: Simplifying training of large, sparse
models, in: ICML, 2021.

\bibitem{mixtral} A.~Q. Jiang, A.~Sablayrolles, A.~Roux, et al., Mixtral of
experts, arXiv:2401.04088 (2024).

\bibitem{olmoe} N.~Muennighoff, L.~Soldaini, D.~Groeneveld, et al.,
{OLMoE}: Open mixture-of-experts language models, arXiv:2409.02060 (2024).

\bibitem{kvv} R.~M. Karp, U.~V. Vazirani, V.~V. Vazirani, An optimal
algorithm for on-line bipartite matching, in: ACM STOC, 1990, pp. 352--358.

\bibitem{adwords} A.~Mehta, A.~Saberi, U.~Vazirani, V.~Vazirani, {AdWords}
and generalized online matching, Journal of the ACM 54 (5) (2007).

\bibitem{buchbinder} N.~Buchbinder, J.~(S.)~Naor, The design of competitive
online algorithms via a primal--dual approach, Foundations and Trends in
Theoretical Computer Science 3 (2--3) (2009) 93--263.

\bibitem{devanurhayes} N.~R. Devanur, T.~P. Hayes, The AdWords problem:
online keyword matching with budgeted bidders under random permutations,
in: ACM EC, 2009, pp. 71--78.

\bibitem{awy} S.~Agrawal, Z.~Wang, Y.~Ye, A dynamic near-optimal algorithm
for online linear programming, Operations Research 62 (4) (2014) 876--890.

\bibitem{krtv} T.~Kesselheim, K.~Radke, A.~T\"onnis, B.~V\"ocking, Primal
beats dual on online packing LPs in the random-order model, in: ACM STOC,
2014, pp. 303--312.

\bibitem{blm} S.~Balseiro, H.~Lu, V.~Mirrokni, The best of many worlds:
dual mirror descent for online allocation problems, Operations Research
71 (1) (2023) 101--119.

\bibitem{twochoices} Y.~Azar, A.~Z. Broder, A.~R. Karlin, E.~Upfal,
Balanced allocations, SIAM Journal on Computing 29 (1) (1999) 180--200.

\bibitem{mitzenmacher} M.~Mitzenmacher, The power of two choices in
randomized load balancing, IEEE Transactions on Parallel and Distributed
Systems 12 (10) (2001) 1094--1104.

\bibitem{sinkhorn} M.~Cuturi, Sinkhorn distances: Lightspeed computation of
optimal transport, in: NeurIPS, 2013, pp. 2292--2300.

\bibitem{wikitext} S.~Merity, C.~Xiong, J.~Bradbury, R.~Socher, Pointer
sentinel mixture models, in: ICLR, 2017.

\end{thebibliography}

\end{document}
